\subsection{The case of a monogenous transcendental extension}

LEMMA 1. Let K be a field, v a valuation on K, Γ its order group, Γ' a totally ordered

group containing \(\Gamma\) and \(\xi\) an element of \(\Gamma'\) . There exists a unique valuation \(w\) on \(\mathbf{K}(\mathbf{X})\) such that, for \(\mathbf{P} = \sum_{j} a_j \mathbf{X}^j (a_j \in \mathbf{K})\) , \(w(\mathbf{P}) = \inf_j(v(a_j) + j\xi)\) .

By Proposition 4 of § 3, no. 2, it suffices to show that the formula

\[
w \left(\sum_ {j} a _ {j} \mathbf {X} ^ {j}\right) = \inf _ {j} (v (a _ {j}) + j \xi)\tag{1}
\]

defines a valuation on the ring K{[}X{]}. As

\[
v \left(a _ {j} + b _ {j}\right) + j \xi \geqslant \inf (v \left(a _ {j}\right), v \left(b _ {j}\right)) + j \xi = \inf (v \left(a _ {j}\right) + j \xi , v \left(b _ {j}\right) + j \xi),
\]

it follows that

\[
w (\mathbf {P} + \mathbf {Q}) \geqslant \inf (w (\mathbf {P}), w (\mathbf {Q}))\tag{2}
\]

for \(\mathbf{P},\mathbf{Q}\) in \(\mathbf{K}[\mathbf{X}]\) , equality holding if \(w(\mathbf{P})\neq w(\mathbf{Q})\) . We show that

\[
w (\mathrm{PQ}) = w (\mathrm{P}) + w (\mathrm{Q})\tag{3}
\]

for \(\mathbf{P} = \sum_{j} a_{j} \mathbf{X}^{j}\) and \(\mathbf{Q} = \sum_{j} b_{j} \mathbf{X}^{j}\) . Let \(i\) (resp. \(k\) ) be the least of the integers \(j\) such that \(v(a_{j}) + j\xi\) (resp. \(v(b_{j}) + j\xi\) ) attains its minimum; let \(\alpha\) (resp. \(\beta\) ) denote this minimum. For \(j, j'\) in \(\mathbf{N}\) ,

\[
w (a _ {j} b _ {j ^ {\prime}} \mathbf {X} ^ {j + j ^ {\prime}}) = v (a _ {j}) + j \xi + v (b _ {j ^ {\prime}}) + j ^ {\prime} \xi \geqslant \alpha + \beta ,
\]

whence \(w(\mathrm{PQ}) \geqslant \alpha + \beta\) by (2). Consider now the term \(c\mathbf{X}^{i + k}\) of degree \(i + k\) in PQ; then \(c = \sum_{n \in \mathbf{Z}} a_{i + n} b_{k - n}\) ; by the choice of \(i\) and \(k\) , the element

\[
w (a _ {i + n} b _ {k - n} \mathbf {X} ^ {i + k}) = v (a _ {i + n}) + (i + n) \xi + v (b _ {k - n}) + (k - n) \xi
\]

takes its minimum value \(\alpha + \beta\) once and once only with \(n = 0\) ; hence \(w(c\mathbf{X}^{i + k}) = \alpha + \beta\) , whence, by (1),

\[
w (\mathrm{PQ}) = \alpha + \beta = w (\mathrm{P}) + w (\mathrm{Q}).
\]

PROPOSITION 1. Let K be a field, v a valuation on K, \(\Gamma\) its order group, \(\Gamma'\) a totally ordered group containing \(\Gamma\) and \(\xi\) an element of \(\Gamma'\) such that the relations \(n\xi\in\Gamma\) , \(n\in Z\) imply n=0. Then there exists a unique valuation w on K(X) with values in \(\Gamma'\) and extending v such that \(w(\mathbf{X})=\xi\) . The residue field of w is equal to that of v and its order group is the subgroup \(\Gamma+Z\xi\) of \(\Gamma'\) .

We show first the uniqueness of \(w\) . Let \(\mathbf{P} = \sum_{j} a_j \mathbf{X}^j\) be an element of \(\mathbf{K}[\mathbf{X}]\) . Then \(w(a_j \mathbf{X}^j) = v(a_j) + j\xi\) , which shows that the monomials \(a_j \mathbf{X}^j\) such that \(a_j \neq 0\) have distinct values for \(w\) . It follows that \(w(\mathbf{P}) = \inf_j (v(a_j) + j\xi)\) , which shows both the uniqueness of \(w\) on \(\mathbf{K}[\mathbf{X}]\) (hence also on \(\mathbf{K}(\mathbf{X})\) ) and the fact that the order group of \(w\) is \(\Gamma + \mathbf{Z}\xi\) . It is further seen that, if \(\mathbf{P} \neq 0\) , we may write \(\mathbf{P} = a\mathbf{X}^n (1 + u)\) , where \(a\in \mathbf{K}^*\) , \(n\in \mathbf{N}\) , \(u\in \mathbf{K}(\mathbf{X})\) and \(w(u) > 0\) ; every element \(\mathbf{R}\neq 0\) of \(\mathbf{K}(\mathbf{X})\) can therefore be written in the form

\[
\mathrm{R} = b \mathrm{X} ^ {n} (1 + u ^ {\prime}),
\]

where \(b \in \mathbf{K}^*\) , \(n \in \mathbf{Z}\) , \(u' \in \mathbf{K}(\mathbf{X})\) and \(w(u') > 0\) ; then \(w(\mathbf{R}) = v(b) + n\xi\) , hence \(w(\mathbf{R}) = 0\) if and only if \(v(b) = 0\) and \(n = 0\) ; thus, when \(w(\mathbf{R}) = 0\) , \(\mathbf{R}\) and \(b\) are congruent modulo the ideal of \(w\) , which shows that the residue field of \(w\) is equal to that of \(v\) .

Finally the existence of w follows from Lemma 1.

PROPOSITION 2. Let K be a field, v a valuation on K, \(\Gamma\) its order group and k its residue field. There exists a unique valuation w on K(X) extending v such that \(w(X) = 0\) and the image t of X in the residue field \(k'\) of w is transcendental over k. The order group of w is equal to that of v and its residue field is \(k(t)\) .

To show the uniqueness of \(w\) , it will suffice for us to show that, if \(\mathbf{P} = \sum_{j} a_j \mathbf{X}^j\) is a non-zero element of \(\mathbf{K}[\mathbf{X}]\) , then

\[
w (\mathbf {P}) = \inf _ {j} (v (a _ {j})).
\]

We may divide P by an element of \(K^{*}\) and suppose that \(v(a_{j}) \geqslant 0\) for all j and that one of the \(v(a_{j})\) is zero. As \(w(\mathbf{X}) = 0\) , P then belongs to the ring of w; writing \(\bar{a}_{j}\) for the canonical image of \(a_{j}\) in k, the canonical image of P in the residue field \(k'\) is \(\sum_{j} \bar{a}_{j} t^{j}\) ; as t is transcendental over k and one of the \(\bar{a}_{j}\) is non-zero, this image is non-zero, whence

\[
w (\mathbf {P}) = 0 = \inf _ {j} (v (a _ {j})).
\]

We now show the existence of \(w\) . The formula \(w(\mathbf{P}) = \inf_{j}(v(a_j))\) (for \(\mathbf{P} = \sum_{j} a_j \mathbf{X}^j\) ) defines a valuation \(w\) on \(\mathbf{K}(\mathbf{X})\) , by virtue of Lemma 1, and \(w\) obviously has the same order group as \(v\) . Then \(w(\mathbf{X}) = 0\) . We show that the canonical image \(t\) of \(\mathbf{X}\) in the residue field \(k'\) of \(w\) is transcendental over \(k\) : if \(\sum_{j} \bar{a}_j t^j = 0\) , where \(\bar{a}_j \in k\) for all \(j\) , then, denoting by \(a_j\) a representative of \(\bar{a}_j\) in the ring of \(v\) , \(w\left(\sum_{j} a_j \mathbf{X}^j\right) > 0\) ; whence \(v(a_j) > 0\) for all \(j\) , hence \(\bar{a}_j = 0\) for all \(j\) . We show finally that \(k' = k(t)\) : every element \(\mathbf{R}\) of \(\mathbf{K}(\mathbf{X})\) may be written \(\mathbf{R} = c\left(\sum_{j} a_j \mathbf{X}^j\right) / \left(\sum_{j} b_j \mathbf{X}^j\right)\) , where \(c, a_j, b_j\) are in \(\mathbf{K}\) , \(v(a_j) \geqslant 0\) and \(v(b_j) \geqslant 0\) for all \(j\) , one of the \(v(a_j)\) and one of the \(v(b_j)\) being zero; then \(w(\mathbf{R}) \geqslant 0\) if and only if \(v(c) \geqslant 0\) ; denoting by \(f\) the canonical homomorphism of the ring of \(w\) onto \(k'\) ,

\[
f (\mathrm{R}) = f (c) \left(\sum_ {j} f (a _ {j}) t ^ {j}\right) / \left(\sum_ {j} f (b _ {j}) t ^ {j}\right),
\]

which proves our assertion.

Remark. It should not be thought that the two types of extensions of v to K(X) which we have just met are the only ones; there may exist a third type of extension, where \(\Gamma'/\Gamma\) is a torsion group and \(k'\) a (not necessarily finite) algebraic extension of k. This third type is not necessarily provided by the procedure described in Lemma 1 (cf.~§3, Exercise 1).

